\begin{abstract}
The EU Circular Material Use Rate (CMUR) is disseminated annually, although its waste-treatment source is observed biennially and intervening years are estimated. Which temporal claims remain supportable when this provenance is respected? We compare the released annual series with released CMUR restricted to admissible source-anchor years for medium-term trends over 2010--2022 across the 27 EU Member States and EU27\_2020, and assess annual change, turning and acceleration through 2022 using source-aware sensitivity trajectories. Under a common two-sided 95\% Student-$t$ ordinary least-squares rule, Member-State direction comparisons yield 12 agreements, no resolved-class disagreements and 15 unresolved cases. Anchor-only model-based intervals are wider in all 27 Member States, with a median paired anchor-only/annual interval-width ratio of 1.858. At a sign-support threshold of 0.95, an unscaled four-year pseudo-gap residual stress test yields 10 annual-change, two turning and one acceleration agreements among 324, 297 and 297 claim instances, respectively. Scaling the same sampled residuals by $1/\sqrt{2}$ in a secondary Brownian-motivated amplitude scenario increases these agreements to 29, 13 and seven; all remaining comparisons are unresolved in both scenarios. These sensitivities leave the true two-year error distribution unidentified. Majority unresolvedness persists for every metric under both scenarios across the tested 0.80--0.99 decision-threshold grid. The evidence identifies a CMUR-specific scale-dependent temporal-support boundary: medium-term direction is comparatively robust where resolved, while model-based precision is reduced and most local temporal claims remain unresolved under the specified sensitivity criteria.
\end{abstract}
